\section{Thermal calculation at inverse temperature \texorpdfstring{$21/2$}{21/2}}
\label{sec:thermal120}

We prove the $b=21/2$ table in Lemma~\ref{lem:thermal-inputs}.
The calculation uses the actual physical Hamiltonians at lengths
$n=6,8,10$, with twists $1,P,C$.  All arithmetic below is integral or
rational; in the cyclic case complex numbers are represented in
$\Z[\omega]$, where $\omega=e^{2\pi i/3}$.

\subsection{The finite matrices and their columns}

Fix $b=21/2$ and one of these lengths, and write $H=H_n(g)$.
After the global rotations in
Section~\ref{sec:model}, the three twists have angles
$\theta=0,\pi,2\pi/3$ about the weight-basis axis.
Set
\begin{equation}\label{eq:thermal120-matrices}
 W=I-\frac{H_n(g)+(3n/2)I}{16},\qquad B=16W.
\end{equation}
Lemma~\ref{lem:energy} gives $\norm W\le1$ and $\norm B\le16$.
In a fixed total-weight sector, the rows and columns are words
$x\in\{-1,0,1\}^n$ with fixed $\sum_jx_j$.  The diagonal of $B$ is
\[
 B_{x,x}=16-\frac{3n}{2}-\sum_{j=0}^{n-1}x_{j-1}x_j.
\]
For each $j$, put $k=j+1\pmod n$.  If
$y_j=x_j-\delta$, $y_k=x_k+\delta$, $\delta\in\{-1,1\}$, and all
letters remain in $\{-1,0,1\}$, the corresponding entry is $-1$, except
at $j=n-1,k=0$, where it is $-e^{i\theta\delta}$.
This is exactly the physical seam phase in~\eqref{eq:seam-phase}.
For the cyclic twist its action on a coordinate $u+v\omega$ uses
\begin{equation}\label{eq:thermal120-phases}
 -\omega(u+v\omega)=v+(v-u)\omega,
 \qquad
 -\omega^{-1}(u+v\omega)=(u-v)+u\omega.
\end{equation}
Thus matrix-vector multiplication requires only integer additions and
multiplications in both coordinates.

Only some columns need to be computed.  To justify the reduction, let
$s(x)=(x_{n-1},x_0,\ldots,x_{n-2})$.  The seam formula gives
\[
 H_{s(x),s(y)}
 =e^{i\theta(y_{n-1}-x_{n-1})}H_{x,y}.
\]
Consequently the monomial unitary
$|x\rangle\mapsto e^{-i\theta x_{n-1}}|s(x)\rangle$ commutes with $H$.
Reversal of the word takes $H$ to its complex conjugate, as does negation
of every letter.  Composing either permutation with complex conjugation
gives an antiunitary symmetry.  Each of these symmetries preserves the
norm of the corresponding column of any real exponential of $H$.

In every sector of nonnegative total weight, select the lexicographically
least word in
each orbit under rotations and reversal.  Give it the orbit size as its
multiplicity, doubled for strictly positive total weight to include the
opposite sector.  The multiplicities then sum to $3^n$ and reproduce the
full Frobenius norm of the exact exponential.  No symmetry of the rounded
columns is required: we compare each computed representative with its
own exact column and repeat that comparison with the same multiplicity.

\subsection{A rational polynomial and its error}

Write
\[
 E=e^{84(W-I)}=e^{-(b/2)(H_n(g)+(3n/2)I)},
 \qquad x_n=(3/2-a)bn.
\]
The desired trace is
\begin{equation}\label{eq:thermal120-frobenius}
 Z_n(b,g)=e^{x_n}\norm{E}_F^2,
\end{equation}
where $\norm{\cdot}_F$ is the Frobenius norm.
For $0\le j\le J=256$, define the rational probabilities
\begin{equation}\label{eq:thermal120-probabilities}
 p_j=\frac{84^j/j!}{\sum_{l=0}^{256}84^l/l!},
 \qquad P(W)=\sum_{j=0}^{256}p_jW^j.
\end{equation}
These are Poisson probabilities of mean $84$, conditioned on values at
most $256$.  For a Poisson variable $V$ with that mean,
\[
 \Pr(V>256)\le\frac{\mathbb E(2^V)}{2^{257}}
 =\frac{e^{84}}{2^{257}}
 <\frac{27^{28}}{2^{257}}<2^{-117}.
\]
The conditioning estimate is the coefficient argument of
Fox--Glynn~\cite[Proposition~1]{FoxGlynn1988}, applied here to powers
of a contraction. Since $\norm{W^j}\le1$, conditioning changes the operator average by at
most twice this probability.  Hence
\begin{equation}\label{eq:thermal120-poisson-error}
 \norm{P(W)-E}<2^{-116}.
\end{equation}

Put $Q=2^{56}$ and $q_j=\lfloor Qp_j\rfloor$.  For a representative
column $r$, start at zero and apply, for $j=256,\ldots,0$,
\begin{equation}\label{eq:thermal120-horner}
 z\longleftarrow\operatorname{fl}(Bz/16)+q_j e_r.
\end{equation}
For real coordinates, $\operatorname{fl}$ is the ordinary floor; for
$u+v\omega$ it floors $u$ and $v$ separately.
The ideal suffix at step $j$ is
$Q\sum_{l=j}^{256}p_lW^{l-j}e_r$, whose norm is at most $Q$.
In a sector of dimension $d_0$, the Euclidean norm of the error from a
coordinatewise floor is less than $2\sqrt{d_0}$, and the coefficient floor adds at most
one.  Contraction of $W$ therefore bounds the error at every stage by
\begin{equation}\label{eq:thermal120-column-error}
 257(2\sqrt{3^n}+1)\le125159.
\end{equation}
The zero coefficients above $q_{172}$ may be omitted: the computed
vector is still zero there, and~\eqref{eq:thermal120-column-error}
continues to count all $257$ possible rounding steps.

The norm of the repeated computed-column list, divided by $Q$ and
multiplied by $e^{x_n/2}$, differs from $\sqrt{Z_n(b,g)}$ by at most
\begin{equation}\label{eq:thermal120-scaled-error}
 e^{x_n/2}\sqrt{3^n}
 \left(\frac{257(2\sqrt{3^n}+1)}Q+2^{-116}\right)
 <4\cdot10^{-7}.
\end{equation}
Indeed $\sqrt{3^n}\le243$, $x_n/2<6$, and $e^{x_n/2}<3^6$.
The resulting rational upper bound
$3^6\,243(125159/2^{56}+2^{-116})$ is less than $4\cdot10^{-7}$.
This comparison includes both polynomial truncation and integer rounding.

\subsection{Exact evaluation by packed integers}

For efficiency, the calculation performs~\eqref{eq:thermal120-horner}
on batches of at most $192$ columns.  A list of signed integer digits
$y_0,\ldots,y_{w-1}$ is stored as
\[
 \operatorname{pack}(y)=\sum_{t=0}^{w-1}y_t2^{64t}.
\]
Integer linear arithmetic preserves this identity before the floor,
even if intermediate carries occur.  Arbitrary-precision integers are
used throughout.

Here is why the floor can also be performed on the packed integer.
For either coefficient of $u+v\omega$,
\[
 |u|,|v|\le\frac2{\sqrt3}|u+v\omega|\le2|u+v\omega|.
\]
By~\eqref{eq:thermal120-column-error}, every computed column has norm
at most $Q+125159$.  Thus every coordinate of $Bz$, in either integer
component, has magnitude at most
\begin{equation}\label{eq:thermal120-lane-bound}
 32(Q+125159)=2305843009217699040<2^{63}.
\end{equation}
Let $s_w=\sum_{t=0}^{w-1}2^{64t}$ and define
\[
 \mathrm{bias}=2^{63}s_w,\qquad
 \mathrm{mask}=(2^{60}-1)s_w,\qquad
 \mathrm{bias16}=2^{59}s_w.
\]
For $Y=\operatorname{pack}(y)$ satisfying~\eqref{eq:thermal120-lane-bound},
the operation
\[
 \texttt{(((Y+bias)>{}>4)\&mask)-bias16}
\]
equals $\operatorname{pack}((\lfloor y_t/16\rfloor)_{t=0}^{w-1})$.
To see this, adding the bias makes every base-$2^{64}$ digit lie in
$[0,2^{64})$.  The global right shift moves four low bits from each
higher lane into the four high bits of the preceding lane.  The mask
deletes precisely those crossing bits, retaining
$\lfloor(y_t+2^{63})/16\rfloor$ in lane $t$.
Subtracting $2^{59}$ then gives $\lfloor y_t/16\rfloor$, also for
negative $y_t$.  This proves by induction that the packed recurrence
is exactly~\eqref{eq:thermal120-horner}.

The final coordinates satisfy the smaller bound
$2(Q+125159)<2^{63}$, so biasing and unpacking recovers every integer
coordinate exactly.  If a final coordinate is $u+v\omega$, its squared
modulus is $u^2-uv+v^2$.  Summing these integers with the representative
multiplicities gives an integer $t_{n,g}$, the squared norm of the full
repeated-column list in $Q$-units.

\subsection{The exact totals and trace enclosures}

Integer evaluation of the recurrence gives the following totals.
The accompanying computation implements the matrix entries,
representatives and packed operations specified above.
\begin{center}
\begin{tabular}{ccl}
\toprule
$n$&$g$&$t_{n,g}$\\
\midrule
6&$1$&93637037510293158533967385845216\\
6&$P$&3629519998654831111446492683663\\
6&$C$&9839536903588146388838689926305\\
8&$1$&4947327904487991279560670127283\\
8&$P$&749999795795292705866802820621\\
8&$C$&1301346729026243689610427021457\\
10&$1$&388818111664728958208623998490\\
10&$P$&120458568770462253326886628639\\
10&$C$&167487042402286114495695779656\\
\bottomrule
\end{tabular}
\end{center}

To enclose the scalar factor, put
\[
 L_n=\sum_{j=0}^{100}\frac{x_n^j}{j!},\qquad
 U_n=L_n+2\frac{x_n^{101}}{101!}.
\]
Since $0<x_n<12$ and successive terms after the first omitted term
have ratio at most $x_n/102<1/2$, we have
$L_n<e^{x_n}<U_n$.
For each center $c_{n,g}$ in the $b=21/2$ table of
Lemma~\ref{lem:thermal-inputs}, exact rational comparison gives
\begin{equation}\label{eq:thermal120-rational-enclosure}
 c_{n,g}-2\cdot10^{-6}
 <\frac{L_n t_{n,g}}{Q^2}
 <\frac{U_n t_{n,g}}{Q^2}
 <c_{n,g}+2\cdot10^{-6}.
\end{equation}
Here the centers are interpreted as real numbers, including their
$10^{-6}$ scale.
These inequalities in particular put every computed scaled norm below
$4$.  By~\eqref{eq:thermal120-scaled-error}, replacing its square by the
true trace incurs error less than
$(8+4\cdot10^{-7})4\cdot10^{-7}$.
Together with~\eqref{eq:thermal120-rational-enclosure}, this yields
\[
 |Z_n(21/2,g)-c_{n,g}|
 <2\cdot10^{-6}+(8+4\cdot10^{-7})4\cdot10^{-7}
 <10^{-5}.
\]
This proves the $b=21/2$ table in Lemma~\ref{lem:thermal-inputs}.
